\begin{foldable} % LLM --> DD
    The progress of large language models depends on large-scale supervised corpora, yet the cost of curating and training on such data compounds across fine-tuning cycles, hyperparameter sweeps, and continual learning.
    The real burden is not any single run but the cumulative expenditure across the full development lifecycle.
    Corpus size, not model size, is often the binding constraint.
    Reducing it is therefore a systemic priority with direct implications for cost, carbon footprint, and data governance.
\end{foldable}

\begin{foldable} % DD --> Text DD
    Dataset distillation (DD), first proposed by~\cite{wang2018dataset}, replaces a large corpus $\mathcal{D}$ with a smaller surrogate $\tilde{\mathcal{D}} \ll \mathcal{D}$ such that a model trained on $\tilde{\mathcal{D}}$ performs comparably to one trained on $\mathcal{D}$.
    Put another way: what is the smallest textbook that teaches the same exam?
    This places DD squarely within data-centric AI, complementing model compression by operating on the data rather than the model itself.
    Early DD methods work in image space, optimising synthetic pixels via gradient or trajectory matching under a bi-level meta-learning objective~\cite{cazenavette2022dataset,kim2022dataset,wang2018dataset,wang2022cafe,zhao2020dataset}.
    Text breaks these assumptions: non-differentiable token decoding blocks gradient-based synthesis, semantic coherence requirements mean outputs must be human-interpretable, and retaining verbatim training samples raises privacy concerns.
    Text dataset distillation~(text DD) therefore calls for a fundamentally different approach.
\end{foldable}

\begin{foldable} % Baseline text DD methods
    We argue that a complete text DD framework must satisfy three properties: (1) knowledge-informed sample weighting, so that selection reflects each sample's contribution to the downstream objective; (2) a principled dataset-level objective that optimises distributional coverage rather than ranking samples independently; and (3) human-readable output that permits inspection and downstream reuse.
    We survey five prior methods \cite{li2021data,sucholutsky2021soft,maekawa2023dataset,maekawa2024dilm,tao2024textual} in \S2.1 and show that each violates at least one of these properties.
\end{foldable}

\begin{foldable} % Our method TAKE and contributions
    We propose TAKE (Trajectory-Aware Knowledge Estimation), a text DD method that satisfies all three properties.
    Influence functions~\cite{koh2017understanding} offer a natural starting point, quantifying each sample's counterfactual effect on the downstream loss.
    However, scoring at a single checkpoint introduces the \textit{hard-sample bias} (HSB): at convergence, noisy samples dominate influence scores while clean and moderate samples --- the most valuable for a robust distilled dataset --- are suppressed.
    TAKE corrects this by integrating influence scores along the full training trajectory into a per-sample knowledge score that captures easy, moderate, and boundary samples while down-weighting noisy ones.
    These scores feed a discrete Optimal Transport (OT) objective, which selects a synthetic corpus that best covers the knowledge-weighted training distribution.
    \begin{itemize}
        \item \textbf{Diagnosis:} We formalise HSB in the text DD setting, quantify it, and establish it as the core failure mode of naive influence-based selection. We further identify the absence of a dataset-level objective as a second structural gap.
        \item \textbf{Method:} TAKE corrects HSB via trajectory-integrated knowledge scores, which serve as source weights in a discrete OT objective aligned with the DD goal --- the first text DD method to jointly address both sample weighting and distributional coverage.
        \item \textbf{Empirical:} TAKE achieves competitive downstream accuracy across five language tasks at extreme compression (0.1\% or 20 samples/class), producing human-readable distilled corpora validated on task performance across diverse backbone families.
    \end{itemize}
\end{foldable}

The remainder of this paper is organised as follows: \S~\ref{sec:02-related} reviews related work; \S~\ref{sec:03-theory} presents theoretical motivation; \S~\ref{sec:04-method} details TAKE and our DD pipeline; \S~\ref{sec:05-experiments} reports experiments and discussion; \S~\ref{sec:06-limits} reflects limitations and ethical concerns; \S~\ref{sec:07-conclusion} concludes.
